\begin{result}{E.2.26}{One integral comparison for all threefold and CM geometric factors}{Proved}
Let \(\mathcal A_{\leq3,\mathrm{CM}}\) consist of all abelian varieties
over \(\Qb\) geometrically isogenous to products of simple factors of
dimension at most three and arbitrary CM abelian varieties. Take the
full integral diagram of E.2.16: every geometric isogeny lattice,
products, algebraic correspondences, tensor and exterior-power
constructions, specified projectors, finite-model Artin module maps,
and all Artin objects. With \(P\), its labelled roots, and the quartic
pair as in E.2.17 and E.2.23, there is a specified
\(\sigma_*\in\Gamma_{F_{\mathrm{cm}}}\) for which
\[
 \sigma_*|_P=(12)(34),\qquad
 \rho_A(\sigma_*)\Lambda_A=\Lambda_{{}^{\sigma_*}A}
       \quad(A\in\mathcal A_{\leq3,\mathrm{CM}}).
 \tag{E.2.26a}
\]
The completed maps in (E.2.26a) are the completions of one integral
tensor comparison on this entire diagram, with Tate map \(1\).
Its Artin module maps give the comparison of the two specified real
embeddings by E.2.3. This statement concerns the plain integral
diagram; the labelled real-Frobenius witness remains E.2.7.
\end{result}
\begin{proof}
Throughout, \(\Lambda_A=H_1(A(\C),\Z)\) and the Tate modules
are covariant homology Tate modules. Cohomology uses inverse duals.
We prove the actual finite-family image and then its integral
intersection while retaining every previously selected matrix.

\paragraph{1. Actual groups and arithmetic Mumford--Tate equality.}
For a geometrically simple non-CM threefold \(A\), put
\(V=\Lambda_A\otimes\Q\), choose a labelled polarization \(\psi\),
and write \(H_A=\operatorname{MT}(A)^{\mathrm{der}}\).
The possibilities, as actual groups acting on \(V\), are
\[
\begin{array}{c|c|c}
 \End^0(A)&\text{coefficient dimension}&H_A\\ \hline
 \Q&\dim_{\Q}V=6&\operatorname{Sp}(V,\psi)\\
 F\text{ totally real cubic}&\dim_FV=2
      &\operatorname{Res}_{F/\Q}\operatorname{SL}_F(V)\\
 E\text{ imaginary quadratic}&\dim_EV=3
      &\operatorname{SU}_E(V,h).
\end{array}
\tag{E.2.26b}
\]
In the second row, the unique \(F\)-alternating form \(\Psi\)
satisfies \(\psi=\operatorname{Tr}_{F/\Q}\Psi\). In the third,
fix \(a\in E^\times\) with \(\bar a=-a\); the Hermitian form
\(h\) is determined by \(\psi=\operatorname{Tr}_{E/\Q}(a h)\).
It has signature \((2,1)\), up to interchange. The remaining simple
threefold case has sextic CM endomorphism field. Moonen--Zarhin's
classification and trace-form identifications give exactly these
groups, and the imaginary-quadratic Hodge group is the full unitary
group \(\operatorname U_E(V,h)\)
\cite[Section 2.3 and Proposition 2.4, author p.~9;
Theorem 2.7, author p.~10]{GenMoonenZarhin1999}.
Consequently its full Mumford--Tate group is
\(\operatorname{GU}_E(V,h)\): the product of the Hodge group and
the rational scalar torus is connected of dimension ten, the
dimension of \(\operatorname{GU}\). The lower-dimensional factors
retain the exact groups and forms of E.2.25.

Every group \(H_A\) in (E.2.26b) is semisimple simply connected.
Its \(\Q\)-simple factors have noncompact real points:
\(\operatorname{Sp}_6(\R)\), \(\operatorname{SL}_2(\R)^3\),
or \(\operatorname{SU}(2,1)\), respectively.

We verify the arithmetic equality before applying an adelic image
theorem. Under the actual Betti--\'{e}tale comparison it is
\[
 G_\ell(A)^0=\operatorname{MT}(A)_{\Q_\ell}.
 \tag{E.2.26c}
\]
Tankeev's original proof in prime dimension different from seven
gives the semisimple Lie dimensions needed here: eight for the
imaginary-quadratic threefold, nine for the real-cubic case, and
twenty-one for \(\End^0(A)=\Q\). The natural inclusion into the
Mumford--Tate Lie algebra is used explicitly in that proof. Its full
unitary equality includes the hypothesis \(c_\ell^0\ne0\), which
we now verify
\cite[Lemma 3 and proof, pp.~532--535; type-I cases,
pp.~535--539]{GenTankeev1988}.

Commelin's exact center theorem, under the same comparison, is
\[
 Z(G_\ell(A)^0)^0
     = Z(\operatorname{MT}(A))^0_{\Q_\ell}.
 \tag{E.2.26d}
\]
Its proof has the following maps. Faltings' endomorphism theorem
identifies the centralizer of the connected arithmetic image with
the geometric endomorphism algebra; Commelin's Lemma 4.15 then
places its connected center inside the Mumford--Tate center. A
faithful representation \(N\) of
\(\operatorname{MT}(A)^{\mathrm{ab}}\) is an abelian CM motive,
by the full Hodge realization theorem and Section 4.16. The CM
Mumford--Tate theorem gives its connected arithmetic image equal
to that torus. Thus the actual projection
\(G_\ell(A)^0\to\operatorname{MT}(A)^{\mathrm{ab}}_{\Q_\ell}\)
is surjective and proves (E.2.26d)
\cite[Theorem 4.13, Lemmas 4.15--4.18 and proof,
pp.~658--659]{GenCommelin2019}.
The source takes a finitely generated definition extension for
\(N\). Its algebraic part over the number field is finite; the
extension is regular over that algebraic part, so restriction of
absolute Galois groups is surjective onto its Galois group.
Consequently the image in the original arithmetic group changes
by finite index and its identity component is unchanged.
In the unitary threefold case (E.2.26d) gives center dimension two,
so \(c_\ell^0\) has dimension one. The qualification in Tankeev's
unitary argument is satisfied. In all three rows of (E.2.26b),
the natural inclusion has equal Lie dimension and hence is equality
of connected groups. This proves (E.2.26c); Lombardo's explicit
prime-dimensional statement supplies a separate source check
\cite[Theorem 2.18, p.~6]{GenLombardoNonsimple2015}.

There is also a direct arithmetic check of the unitary center.
After a model extension defining the \(E\)-action and containing
\(E\), Fit\'e--Goodman construct a CM elliptic curve \(B\) with
\[
 \left(\bigwedge_{E_\lambda}^3 V_\lambda(A)
       \oplus\bigwedge_{E_\lambda}^3\overline{V_\lambda(A)}\right)(-1)
       \simeq V_\ell(B)\otimes_{\Q_\ell}E_\lambda.
 \tag{E.2.26e}
\]
Here the wedge is the determinant over the coefficient field
\(E_\lambda\); the bar denotes the coefficient-conjugate
representation defined in that source, including the isomorphism
between the conjugate coefficient completions
\cite[Theorem 1.1(i), p.~2, and Theorem 3.1 with proof,
pp.~9--10]{GenFiteGoodman2026}.
The determinant characters are therefore the two CM elliptic
characters multiplied by the common cyclotomic character.
The torus map is
\(t\mapsto t\operatorname{Norm}_{E/\Q}(t)\); on split
characters its matrix is
\(\left(\begin{smallmatrix}2&1\\1&2\end{smallmatrix}\right)\),
of determinant three. It is an isogeny of a two-dimensional torus,
which verifies the center dimension directly. This use of
(E.2.26e) is an equality of the stated arithmetic representations.

For every finite family of these factors, arithmetic
Mumford--Tate equality holds for their product by Commelin's
product theorem
\cite[Theorem 10.5, p.~676]{GenCommelin2019}.
We use its actual joint group, whose shared factors are handled
next.

\paragraph{2. The joint derived group and an exact rational splitting.}
For a finite family let
\(M=\operatorname{MT}(\bigoplus_iV_i)\) and \(H=M^{\mathrm{der}}\).
The projections \(M\to M_i\) and \(H\to H_i\) are surjective,
and the faithful sum representation embeds \(H\) as a closed
connected subgroup of \(\prod_iH_i\).

Here is an explicit proof that \(H\) is simply connected.
Over \(\Qb\), write its universal cover as
\(H^{\mathrm{sc}}=\prod_jS_j\). For an absolutely simple
factor \(T\) of a target \(H_i\), the images of the \(S_j\)
in \(T\) are commuting connected normal subgroups. Surjectivity
implies that exactly one is nontrivial. Its map \(S_j\to T\)
is a central isogeny, hence an isomorphism because \(T\) is
simply connected. Every \(S_j\) reaches at least one such
target, since otherwise the faithful embedding of \(H\) would
kill its positive-dimensional image. An element of the kernel
of \(H^{\mathrm{sc}}\to H\) therefore has trivial coordinate
in every \(S_j\), using its isomorphic target projection.
The kernel is trivial, and this descends to \(\Q\).

Every \(\Q\)-simple factor \(S\) of \(H\) maps nontrivially
to some \(\Q\)-simple \(H_i\). The image is connected normal,
so it is all of \(H_i\); \(\Q\)-simplicity gives finite
kernel, and the simply connected target makes the map an
isomorphism. Every such \(S\) therefore has noncompact real
points and is one of the forms in (E.2.26b) or E.2.25.

For the joint derived groups of the first \(n\) generators,
the forgetting map \(p_n:H_n\to H_{n-1}\) is surjective.
Over \(\Qb\) it discards a subset of simple factors and maps
the remaining ones isomorphically onto the target factors.
The kernel subset and its complement are Galois stable, giving
the actual \(\Q\)-decomposition
\[
 H_n=L_n\times C_n,\qquad L_n=\ker p_n,
 \qquad p_n|_{C_n}:C_n\xrightarrow{\sim}H_{n-1}.
 \tag{E.2.26f}
\]
Let \(j_n:H_{n-1}\to H_n\) be its inverse section.
The group \(L_n\) is connected semisimple simply connected,
and all its \(\Q\)-simple factors have noncompact real points.
The possibility \(L_n=1\) is included. In particular (E.2.26f)
retains all shared diagonal factors in the precise way needed
to keep earlier matrices fixed.

\paragraph{3. Actual prime-supported derived subgroups.}
For a finite product choose a number field \(k\) defining
the models and endomorphisms, enlarged so that every
\(G_\ell\) is connected. The common connectedness field is
the one recalled in Cadoret--Moonen's Section 1.1.
Their Theorem A(i), with Section 1.3, applies to (E.2.26c):
each local image is open in \(M(\Z_\ell)\), and for all
but finitely many \(\ell\) it contains
\([M(\Z_\ell),M(\Z_\ell)]\)
\cite[Theorem A(i), Sections 1.1--1.3]{GenCadoretMoonen2020}.
The integral models are the closures in the actual Betti
lattice; outside a finite set they are the smooth reductive
models of the groups above.

Serre's independence theorem, applied to the full family of
Tate modules of this product, gives a finite extension \(k'/k\)
with the actual equality
\[
 \rho(\Gamma_{k'})=\prod_\ell\rho_\ell(\Gamma_{k'}).
 \tag{E.2.26g}
\]
This is the independence defined on p.~542 and proved for
these Tate modules in Theorem 1 and Section 3.1
\cite[pp.~542--545]{GenSerreIndependence2013}.
Apply the integral Mumford--Tate statement again over \(k'\).
For a further finite Galois field \(K/\Q\), the image of
\(\Gamma_K\) is a closed finite-index subgroup of the compact
product (E.2.26g), hence is open. It contains a product of open
local subgroups equal to the full local image outside a finite
set. At a good prime \(\ell\geq7\), the perfectness proved
below gives
\[
 H(\Z_\ell)=[H(\Z_\ell),H(\Z_\ell)]
       \subset[M(\Z_\ell),M(\Z_\ell)].
\]
At any remaining prime, intersect a local open subgroup with
\(H(\Q_\ell)\). We obtain the actual containment
\[
 W\subset\rho(\Gamma_K),\qquad
 W\subset H(\mathbf A_f)\text{ compact open},
 \quad W_\ell=H(\Z_\ell)\text{ for almost all }\ell.
 \tag{E.2.26h}
\]
Every element of \(W\) supported at one prime and one local
simple factor has an actual Galois preimage.

\paragraph{4. Finite groups and the full integral unitary lifting.}
The good reductions are
\(\operatorname{SL}_2(\mathbf F_q)\),
\(\operatorname{Sp}_4(\mathbf F_\ell)\),
\(\operatorname{Sp}_6(\mathbf F_\ell)\),
\(\operatorname{SL}_3(\mathbf F_\ell)\), or
\(\operatorname{SU}_3(\mathbf F_\ell)\).
For \(\ell\geq7\) they are perfect and their projective
quotients are simple. The opened matrix proofs are Taylor's
Theorems 4.4--4.5, 8.7--8.8, and 10.15; the last explicitly
proves unitary-transvection generation and perfectness
\cite[pp.~22--23, 72--74, 122--123]{GenTaylorClassical2021}.
A normal subgroup containing a noncentral element surjects
onto the simple projective quotient. The remaining quotient
of the full group is then a quotient of its center, hence
abelian, and perfectness makes it trivial.

We prove the integral lifting statement for \(\operatorname{SU}_3\).
Let \(\mathcal O/\Z_\ell\) be unramified quadratic, choose
a trace-zero unit \(\vartheta\), and use
\[
 J=\begin{pmatrix}0&0&1\\0&1&0\\1&0&0\end{pmatrix},\qquad
 Q(a,b)=\begin{pmatrix}1&-\bar a&b\\0&1&a\\0&0&1\end{pmatrix},
 \quad a=x+y\vartheta,\quad
 b=-a\bar a/2+z\vartheta.
 \tag{E.2.26i}
\]
The exact identities are
\(\det Q=1\), \(\bar Q^{\,t}JQ=J\), and
\(Q(a,b)^{-1}=Q(-a,\bar b)\). Thus both \(Q\) and
its inverse have polynomial degree at most two, and
\(\operatorname{Ad}Q\) has degree at most four.
The matrix \(u=1+N\), \(N=\vartheta E_{13}\), is an
integral unitary transvection with \(N^2=0\).

Suppose \(R\) is a closed subgroup of
\(\operatorname{SU}_3(\Z_\ell)\) surjecting modulo \(\ell\).
Choose \(g\in R\) lifting \(u\bmod\ell\), and write
\(g=u(1+\ell B)\bmod\ell^2\). Then
\[
 g^\ell=u^\ell\left(1+\ell
       \sum_{j=0}^{\ell-1}\operatorname{Ad}(u^{-j})B\right)
       \pmod{\ell^2}.
\]
Here
\(\operatorname{Ad}(u^{-j})B
 =B+j(BN-NB)-j^2NBN\).
The sums of \(1,j,j^2\) vanish modulo \(\ell\geq5\),
so \(g^\ell=1+\ell N\bmod\ell^2\).
The first congruence layer of \(R\) contains \(N\), and
the span of its conjugates is stable under the finite unitary
group. For \(\ell\geq7\), polynomial interpolation in
\(x,y,z\) in (E.2.26i) makes this subspace stable under
bracketing with the six root derivatives
\[
\begin{aligned}
 P_x&=-E_{12}+E_{23},&
 P_y&=\vartheta(E_{12}+E_{23}),&P_z&=\vartheta E_{13},\\
 M_x&=E_{21}-E_{32},&
 M_y&=\vartheta(E_{21}+E_{32}),&M_z&=\vartheta E_{31}.
\end{aligned}
\]
Their brackets
\[
 [P_x,M_x]=\operatorname{diag}(-1,0,1),\qquad
 [P_x,M_y]=\vartheta\operatorname{diag}(-1,2,-1)
\]
together with the six root matrices form a basis of
\(\mathfrak{su}_3(\mathbf F_\ell)\).
The conjugate span is therefore a nonzero Lie ideal.
After quadratic scalar extension this Lie algebra is
\(\mathfrak{sl}_3\), simple in characteristic greater
than three, so the first layer is full. Smoothness identifies
all successive congruence layers with that Lie algebra, and
\[
 (1+\ell^mX)^\ell
      =1+\ell^{m+1}X\pmod{\ell^{m+2}},\qquad m\geq1,
 \tag{E.2.26j}
\]
gives each next layer. Hence \(R\) surjects modulo every
\(\ell^m\), and closedness gives the full integral group.

For \(\operatorname{SL}_3\), use the elementary roots
\(1+tE_{ij}\); their conjugate logarithms span
\(\mathfrak{sl}_3\), and the same calculation applies.
For \(\operatorname{SL}_2\) over an unramified local
field \(F_v\), use transvections with every parameter in
its residue field \(\mathbf F_q\). The first layer
contains all residue multiples of \(E_{12}\) and their
conjugates, hence \(\mathfrak{sl}_2(\mathbf F_q)\) as
an \(\mathbf F_\ell\)-space; (E.2.26j) gives all layers.
For \(\operatorname{Sp}_6\), the logarithms
\(v\mapsto\psi(v,u)u\) span the Lie algebra: under
\(\mathfrak{sp}(V)\simeq\operatorname{Sym}^2(V)\)
they are the squares \(u^2\), which span because
\(2uv=(u+v)^2-u^2-v^2\). The same power calculation
gives all layers, including the already used
\(\operatorname{Sp}_4\) case.
The closed commutator subgroup of each good integral group
surjects modulo \(\ell\) by finite perfectness; the lifting
proof just given makes it the whole integral group. This
establishes the integral perfectness used in (E.2.26h).

\paragraph{5. Complex conjugation on every actual factor.}
Enlarge the finite set of geometric isogeny classes to its
finite \(\Gamma_\Q\)-orbit. For the fixed actual complex
conjugation \(c\), choose isogenies
\(\phi_i:{}^c A_i\to A_{\pi(i)}\), and let
\(C_i=\phi_{i,*}c_i\) on rational Betti homology.
Their block sum \(C\) gives the actual intertwining map;
choose a finite Galois \(K/\Q\) defining every map.
For \(\gamma\in\Gamma_K\), naturality gives
\[
 C_i\rho_i(c\gamma c)C_i^{-1}
         =\rho_{\pi(i)}(\gamma),\qquad
 C h(z)C^{-1}=h(\bar z).
 \tag{E.2.26k}
\]
Thus \(\theta=\operatorname{Ad}(C)\) acts on \(M\)
and \(H\). Its square is trivial there because \(C^2\)
is block diagonal with algebraic endomorphisms, which
centralize \(M\).

On a \(\theta\)-fixed \(\Q\)-simple factor \(S\)
of \(H\), this automorphism cannot be the identity.
Indeed the actual quotient \(M\to S^{\mathrm{ad}}\)
would turn (E.2.26k) into
\(h_{\mathrm{ad}}(z)=h_{\mathrm{ad}}(\bar z)\).
On the unit circle this gives
\(h_{\mathrm{ad}}(z)=h_{\mathrm{ad}}(z)^{-1}\), so
connectedness makes the homomorphism trivial there.
It is also trivial on the real weight torus, since
weight is central. The resulting trivial adjoint
Hodge homomorphism contradicts the fact that this
nontrivial quotient of \(M\) is generated as a
\(\Q\)-group by its image.

Write \(S=\operatorname{Res}_{F/\Q}T\), with \(T\)
absolutely simple over its centroid \(F\). The action
permutes the places of \(F\) over \(\ell\). On a
fixed place, a nontrivial action on \(F\) stays
nontrivial on the centroid \(F_v\). When it fixes
\(F\), identity after the faithful scalar extension
\(F\to F_v\) would force identity over \(F\).
Thus each fixed local simple factor has nontrivial
automorphism; the other local factors occur in exchanged
pairs. Away from finitely many primes all these maps
and models are integral smooth and the fixed-factor
Lie automorphisms remain nontrivial after reduction.

For such a finite factor some
\(g\theta(g)^{-1}\) is noncentral. If every one were
central, this assignment would be a homomorphism from
a perfect group to its center, hence trivial. Then
\(\theta\) would fix every finite point. Interpolation
in the root matrices above, of degree at most two,
would make its differential fix their Lie generators,
contradicting the nontrivial reduced differential.
When the centroid automorphism is nontrivial, using
every residue parameter in those root matrices
would also force that residue-field automorphism
to be trivial. This proves the required noncentral
commutator in each case.

\paragraph{6. Actual supported commutators over \(F_{\mathrm{cm}}\).}
Put \(U_0=\Gamma_K\cap\Gamma_{\Q^{\mathrm{cm}}}\) and
\(N=\rho(U_0)\). This is a compact normal subgroup
of \(\rho(\Gamma_K)\). Actual \(c\) is central on
\(\Q^{\mathrm{cm}}\), so
\[
 [\gamma,c]\in U_0,
 \qquad g\theta(g)^{-1}\in N
       \quad(\gamma\in\Gamma_K, g=\rho(\gamma)).
 \tag{E.2.26l}
\]
We construct subgroups of \(H(\mathbf A_f)\)
inside \(N\); this accommodates any finite action
on the joint central CM torus.

At a good prime, the independently supported full
factors in (E.2.26h), the noncentral commutator in
(E.2.26l), normality, finite projective simplicity,
and integral lifting prove that \(N\) contains
the whole supported fixed factor. For a pair of
factors exchanged by \(\theta\), the first
commutator has two coordinates. A second commutator
with an element supported in the first factor has
only that coordinate. Perfectness makes the closed
group generated by these second commutators the
full factor, and the second factor follows likewise.

At an exceptional prime, use a supported open
subgroup supplied by (E.2.26h). On a fixed local
factor the actual curves
\[
 t\longmapsto\exp(tX)\theta(\exp(tX))^{-1}
\]
have tangent vectors \((1-\theta)X\). For small
\(t\) the matrix exponentials converge and lie
in the actual group: they preserve its defining
bilinear or Hermitian form and have determinant
one for trace-zero \(X\). For a quaternion factor
the same assertions follow from reduced trace
and norm after scalar extension to a splitting
field. A nonzero such tangent and its conjugates
under a local open subgroup span a nonzero ideal
of the local simple Lie algebra, hence all of it.
Choose finitely many resulting actual curves
whose tangents form a basis. Their product map
has invertible differential; the explicit
\(p\)-adic contraction argument of E.2.25 gives
an actual neighborhood of the identity in \(N\)
on this supported factor. For exchanged factors,
second commutators have tangents
\((1-\operatorname{Ad}h)X\). The zero center of
the simple Lie algebra makes their conjugate
span full, giving the same conclusion.

All these constructions have individual prime
and factor support. Finite products and closedness
therefore place a compact open subgroup of
\(H(\mathbf A_f)\) inside \(N\).
The closed derived subgroup of \(U_0\) fixes
every \(E^{\mathrm{ab}}\) for CM \(E\), because
\(U_0\) fixes \(E\) and its restriction to
the abelian Galois group kills commutators.
Compactness gives the exact image equality
\[
 \rho\bigl(\overline{[U_0,U_0]}\bigr)
          =\overline{[N,N]}.
 \tag{E.2.26m}
\]
The right side contains all good integral factors
by their proved perfectness. At an exceptional
prime the closed derived subgroup of a local
open subgroup is open: its Lie algebra is
perfect, and a finite product of the actual
commutator curves with spanning tangents gives
the same local contraction proof. Consequently
\[
 \rho(\Gamma_K\cap\Gamma_{F_{\mathrm{cm}}})
       \text{ contains a compact open subgroup of }
       H(\mathbf A_f).
 \tag{E.2.26n}
\]
When the finite family was enlarged to its Galois
orbit, project back to the original joint derived
group. The projection splits over \(\Q\) by
(E.2.26f), so on finite adeles it is an open
direct-product projection. Its compact open
image still has actual preimages in the same
Galois group. Thus (E.2.26n) holds for every
original finite family as well.

\paragraph{7. Integral approximation with every earlier matrix fixed.}
Let \(\sigma\in\Gamma_{F_{\mathrm{cm}}}\) have
fixed conjugate targets for the first \(n\)
generators. E.2.11 supplies a rational tensor
comparison \(T_0\) on the full rational abelian
category with prescribed all-CM comparison
\(\alpha(\sigma)\), Artin action, and Tate map
one. On the joint representation put
\[
 a=T_0^{-1}\rho(\sigma)\in H_n(\mathbf A_f).
 \tag{E.2.26o}
\]
Here the group membership follows from an actual
restriction map. Let \(D\) be the kernel of
the full abelian-to-CM Tannaka restriction.
Its image in \(M_n\) is connected semisimple,
by the exact kernel theorem used in E.2.11,
and the quotient of \(M_n\) by this image is
a quotient of the geometric CM protorus.
It is a torus, so the image contains
\(M_n^{\mathrm{der}}\); semisimplicity gives
the reverse inclusion. Its image is exactly
\(H_n\). The difference fixes the CM point
at every prime, hence its finite projection
belongs to \(H_n(\Q_\ell)\). Its entries
and inverse are integral outside the finitely
many denominators of \(T_0\), proving
(E.2.26o).

Suppose the old integral matrices are fixed as
\(S\). Since \(\rho(\sigma)\) has exactly
those matrices, \(p_n(a)=T_{0,\mathrm{old}}^{-1}S\)
is rational. It satisfies the defining
\(\Q\)-equations of \(H_{n-1}\) in every
completion, so lies in \(H_{n-1}(\Q)\).
Set
\[
 b=j_n(p_n(a))\in H_n(\Q),\qquad
 b^{-1}a\in L_n(\mathbf A_f).
 \tag{E.2.26p}
\]
Freeze the current restriction to a finite
Galois model field \(K_n\). The exact
stabilizer of this cylinder and all old
matrices is
\[
 \mathcal K_n=\Gamma_{F_{\mathrm{cm}}}
       \cap\Gamma_{K_n}\cap\ker(\rho_{\mathrm{old}}).
 \tag{E.2.26q}
\]
Intersect the compact open subgroup in
(E.2.26n) with \(\ker p_n=L_n(\mathbf A_f)\).
Every element in that intersection has an
actual preimage fixing all old Tate modules,
hence a preimage in \(\mathcal K_n\).
It supplies a compact open subgroup
\(U\subset L_n(\mathbf A_f)\) of its image;
shrink it into the integral lattice stabilizer.

D.4.1 applies to \(L_n\), whose simple factors
are simply connected and have noncompact real
points. Thus \(L_n(\Q)\) is dense in
\(L_n(\mathbf A_f)\). Choose
\[
 l\in L_n(\Q)\cap b^{-1}aU,
 \qquad h=bl\in H_n(\Q).
\]
Then \(p_n(h)=p_n(a)\) and \(h=au\)
for some \(u\in U\). The rational comparison
\(T=T_0h\) retains all old matrices and
its completion is \(\rho(\sigma)u\).
Since \(u\) has an actual preimage
\(k\in\mathcal K_n\), that completion
is \(\rho(\sigma k)\). It is an integral
isomorphism at every prime from the source
lattice onto the fixed target lattice.
Both \(T\) and \(T^{-1}\) have rational
entries in \(\widehat\Z\), and
\(\Q\cap\widehat\Z=\Z\) makes them
integer matrices. For \(L_n=1\),
(E.2.26p) already makes \(a\) rational
and the same conclusion holds with \(u=1\).
This proves the exact induction step.

\paragraph{8. A specified simultaneous element and every geometric lattice.}
Enumerate the geometric isogeny classes of
simple non-CM abelian varieties of dimension
at most three, using the fixed algebraic-data
ordering of E.2.25 to choose their models
and polarizations. Use its integral-basis
rule on each model and conjugate: a chosen
algebraic basis of regular differentials
embeds homology as a full lattice in
\(\C^g\); minimize the sum of squared
period norms of an ordered integral basis
and then its real-coordinate lexicographic
tuple. A first available basis supplies a
finite norm bound, and the full lattice
has finitely many vectors below that bound.
Thus the minimizing basis and finite
tie-break exist. This specifies choices
without a computability assertion.

Start from the nonempty compact fibre
\[
 C_0=\{\sigma\in\Gamma_{F_{\mathrm{cm}}}:
                  \sigma|_P=(12)(34)\}.
\]
Let \(K_n\) be the fixed increasing finite
Galois exhaustion of \(\Qb\), enlarged
at step \(n\) to define the selected
models, conjugate targets, and maps above.
In the current fibre choose the first
finite \(K_n\)-cylinder that meets it;
this freezes the conjugate targets. The
preceding step gives at least one integral
next matrix in that cylinder with all old
matrices unchanged. Choose the first such
integer matrix in entry-height and
lexicographic order. Its fibre \(C_n\)
is compact and nonempty. These fibres are
nested. Their intersection is nonempty,
and contains exactly one \(\sigma_*\),
because its specified restrictions to
the \(K_n\) exhaust \(\Qb\).
This is a specified actual Galois element
with the required quartic action and
integral transports on every enumerated
generator.

Every CM abelian variety has the actual
algebraic isomorphism \(\alpha_A\) of
E.2.6 with completion \(\rho_A(\sigma_*)\).
Products carry direct sums. For any
geometric isogeny \(q:B\to A\) from
such a product, put
\[
 L_A=q_*^{-1}\Lambda_A\subset H_1(B(\C),\Q),
 \qquad \Lambda_B\subset L_A\subset\tfrac1m\Lambda_B,
 \quad L_A/\Lambda_B=\ker(q)\subset B[m]
\]
for a suitable \(m\). Actual Galois
transport on the finite kernel agrees
with the selected integral map on \(B\)
modulo \(m\). That rational map
therefore carries \(L_A\) exactly onto
the lattice for \({}^{\sigma_*}q\).
It induces
\[
 T_A=({}^{\sigma_*}q)_*T_Bq_*^{-1}:
     \Lambda_A\xrightarrow{\sim}
     \Lambda_{{}^{\sigma_*}A},
 \tag{E.2.26r}
\]
whose completion is the actual
\'{e}tale transport. This includes every
specified geometric isogeny lattice.

Compare the actual completed equations
for two choices of \(q\), every product,
isogeny, algebraic correspondence, tensor,
exterior power, and specified projector.
Injectivity of rational matrices into
their completions gives independence
and each required compatibility over
\(\Z\), as in E.2.16. The same argument
preserves the labelled polarizations
with the one Tate line fixed. The
Artin maps are the actual permutations
\(x\mapsto\sigma_*x\). On every finite
model the primitive Artin idempotents
select its conjugate blocks, and their
module maps give exactly the descent
of E.2.3. This proves (E.2.26a) on
the full stated diagram.

The remaining extension requires actual
integral transports for simple non-CM
factors of dimension at least four and
then the further objects and morphisms
of the full pure category. The original
full-category place-separation gate
remains open. E.2.22 is still its exact
all-abelian cover-tower test for this
same \(\sigma_*\).

The finite certificate
\path{code/verify_threefold_groups.py}
checks (E.2.26i), the degree-four
adjoint polynomial, the six root and
two Cartan matrices with an exact
minor \(-1\), the degree-two
transvection power calculation, the
rank-twenty-one symplectic transvection
span, and the determinant-three CM
torus map. The arithmetic image,
supported products, approximation,
and countable construction have the
proofs above. Source access is recorded
in the bibliography and work proof:
Tankeev's case proof was read in the
original indexed PDF, while the missing
Theorem 1 display was not claimed as
visually verified.
\end{proof}
